\documentclass[12pt, english]{article}
%%%%%%%%%%%%%%%%%%%%%%%%%%%%%%%%%%%%%%%%%%%%%%%%%%%%%%%%

% margin %%%%%%%%%%%%%%%%%%%%%%
\usepackage{geometry}
 \geometry{
 a4paper,
 left=1in,
 top=0.9in,
 right=1in,
 bottom=2in,
 }
%%%%%%%%%%%%%%%%%%%%%%%%%%%%%%%%%%%%%%%%%%%%%%%%%%%%%%%%

% page no. %%%%%%%%%%%%%%%%%%%%%%
\usepackage{fancyhdr}
\pagestyle{fancy}
\fancyhf{}
\fancyfoot[C]{\textit{\thepage}}
\setlength{\headheight}{65pt}
%%%%%%%%%%%%%%%%%%%%%%%%%%%%%%%%%%%%%%%%%%%%%%%%%%%%%%%%

% idk-needed %%%%%%%%%%%%%%%%%%%%%%
\usepackage{amsmath,amsthm,amssymb,amsfonts, enumitem, color, comment, graphicx, environ, lipsum}
%%%%%%%%%%%%%%%%%%%%%%%%%%%%%%%%%%%%%%%%%%%%%%%%%%%%%%%%

% code and color %%%%%%%%%%%%%%%%%%%%%%
\usepackage{spverbatim}
\usepackage{xcolor}
\usepackage{listings}
\definecolor{bluegrey}{RGB}{0, 0, 200}
\lstset{
        framexleftmargin=20pt,
        frame=shadowbox,
        rulesepcolor=\color{gray},
        xleftmargin=50pt,
        xrightmargin=30pt,
        breaklines=true,
        basicstyle=\ttfamily,
        showstringspaces=false,
        numbers=left,
        numberstyle=\tiny,}
%%%%%%%%%%%%%%%%%%%%%%%%%%%%%%%%%%%%%%%%%%%%%%%%%%%%%%%%

% heading %%%%%%%%%%%%%%%%%%%%%%
\renewcommand\thesection{\color{blue}\arabic{section}}
\newcommand{\mysection}[2]{\setcounter{section}{#1}\addtocounter{section}{-1}\section{#2}}
%%%%%%%%%%%%%%%%%%%%%%%%%%%%%%%%%%%%%%%%%%%%%%%%%%%%%%%%

% both side indent %%%%%%%%%%%%%%%%%%%%%%
\newenvironment{myindent}
{\par\leftskip25pt\relax\rightskip53.4cm\relax}
{\par\leftskip0cm\relax\rightskip0cm\relax}
%%%%%%%%%%%%%%%%%%%%%%%%%%%%%%%%%%%%%%%%%%%%%%%%%%%%%%%%

% type ^ %%%%%%%%%%%%%%%%%%%%%%
\newcommand{\power}{$^\wedge$}
%%%%%%%%%%%%%%%%%%%%%%%%%%%%%%%%%%%%%%%%%%%%%%%%%%%%%%%%

% document info %%%%%%%%%%%%%%%%%%%%%%
\lhead{Arav Sarma \\ 15193845}  %replace with your name
\rhead{Homework \textcolor{blue}{2} \\ 2025 Speech Processing}
%%%%%%%%%%%%%%%%%%%%%%%%%%%%%%%%%%%%%%%%%%%%%%%%%%%%%%%%

% begin %%%%%%%%%%%%%%%%%%%%%%
\begin{document}
%wwwwwwwwwwwwwwwwwwwwwwwwwwwwwwwwwwwwwwwwwwwwwwwwwwwwwwwww

% section _______________________________________________
\mysection{2}{Programming: expressions, variables}
\vspace{10pt}
% -__-__-__-__-__-__-__-__-__-__-__-__-__-__-__-__-__-__-

% //////////////////////////////////////////////////////
\subsection*{\emph{1.} Print whole numbers 1 to 10} 

\begin{lstlisting}
writeInfoLine: ""
num = 0
for n 1 to 10
    num = num + 1
    appendInfoLine: num
endfor
\end{lstlisting}

\begin{lstlisting}[rulesepcolor=\color{blue}]
writeInfo: ""
for n to 10
    appendInfoLine: n
endfor
\end{lstlisting}

\vspace{10pt}
% //////////////////////////////////////////////////////

% \\\\\\\\\\\\\\\\\\\\\\\\\\\\\\\\\\\\\\\\\\\\\\\\\\\\\\\\\\
\subsection*{\emph{2.} Different name for loop variable}

\begin{lstlisting}
writeInfoLine: ""
num = 0
for baseten 1 to 10
    num = num + 1
    appendInfoLine: num
endfor
\end{lstlisting}

\begin{lstlisting}[rulesepcolor=\color{blue}]
writeInfo: ""
for i to 10
    appendInfoLine: i
endfor
\end{lstlisting}

\vspace{10pt}
% \\\\\\\\\\\\\\\\\\\\\\\\\\\\\\\\\\\\\\\\\\\\\\\\\\\\\\\\\\

% \\\\\\\\\\\\\\\\\\\\\\\\\\\\\\\\\\\\\\\\\\\\\\\\\\\\\\\\\\
\subsection*{\emph{3.} Print whole numbers 10 to 15}

\begin{lstlisting}
writeInfoLine: ""
num = 9
for n 1 to 6
    num = num + 1
    appendInfoLine: num
endfor
\end{lstlisting}

\begin{lstlisting}[rulesepcolor=\color{blue}]
writeInfo: ""
for n from 10 to 15
    appendInfoLine: n
endfor
\end{lstlisting}

\vspace{10pt}
% \\\\\\\\\\\\\\\\\\\\\\\\\\\\\\\\\\\\\\\\\\\\\\\\\\\\\\\\\\

% \\\\\\\\\\\\\\\\\\\\\\\\\\\\\\\\\\\\\\\\\\\\\\\\\\\\\\\\\\
\subsection*{\emph{4.} Even numbers 2 to 16}

\begin{lstlisting}
writeInfoLine: ""
num = 0
for n 1 to 8
    num = num + 1
    appendInfoLine: num * 2
endfor
\end{lstlisting}

\begin{lstlisting}[rulesepcolor=\color{blue}]
writeInfo: ""
for i to 8
    appendInfoLine: i * 2
endfor
\end{lstlisting}

\vspace{10pt}
% \\\\\\\\\\\\\\\\\\\\\\\\\\\\\\\\\\\\\\\\\\\\\\\\\\\\\\\\\\

% \\\\\\\\\\\\\\\\\\\\\\\\\\\\\\\\\\\\\\\\\\\\\\\\\\\\\\\\\\
\subsection*{\emph{5.} Odd numbers 1 to 15}

\begin{lstlisting}
writeInfoLine: ""
num = 0
for n 1 to 8
    num = num + 1
    appendInfoLine: num * 2 - 1
endfor
\end{lstlisting}

\begin{lstlisting}[rulesepcolor=\color{blue}]
writeInfo: ""
for n from 1 to 8
    appendInfoLine: n * 2 - 1
endfor
\end{lstlisting}

\vspace{10pt}
% \\\\\\\\\\\\\\\\\\\\\\\\\\\\\\\\\\\\\\\\\\\\\\\\\\\\\\\\\\

% \\\\\\\\\\\\\\\\\\\\\\\\\\\\\\\\\\\\\\\\\\\\\\\\\\\\\\\\\\
\subsection*{\emph{6.} Whole numbers 1 to n}

\begin{lstlisting}
writeInfoLine: ""
form
    positive number
endform
for a 1 to number
    num = num + 1
    appendInfoLine: num
endfor
\end{lstlisting}

\begin{lstlisting}[title=\textbf{Or:}]
writeInfoLine: ""
form
    integer number
endform
if number < 0
    appendInfoLine: "number is not >=0"
elsif number = 0
    appendInfoLine: 0
    appendInfoLine: 1
elsif number = 1
    appendInfoLine: 1
else
    num = 0
    for n 1 to number
        num = num + 1
        appendInfoLine: num
    endfor
endif
\end{lstlisting}

\begin{lstlisting}[rulesepcolor=\color{blue}]
writeInfo: ""
form
    natural number
endform
for a to number
    appendInfoLine: a
endfor
\end{lstlisting}

\vspace{10pt}
% \\\\\\\\\\\\\\\\\\\\\\\\\\\\\\\\\\\\\\\\\\\\\\\\\\\\\\\\\\

% \\\\\\\\\\\\\\\\\\\\\\\\\\\\\\\\\\\\\\\\\\\\\\\\\\\\\\\\\\
\subsection*{\emph{7.} Modulo Operator}
\vspace{10pt}

\begin{enumerate}
    \item[a.] 0, 1, 0, 1
    \item[b.] When x mod y is 0, y is a factor of x. 
\end{enumerate}
\vspace{10pt}
% \\\\\\\\\\\\\\\\\\\\\\\\\\\\\\\\\\\\\\\\\\\\\\\\\\\\\\\\\\

% \\\\\\\\\\\\\\\\\\\\\\\\\\\\\\\\\\\\\\\\\\\\\\\\\\\\\\\\\\
\subsection*{\emph{8.} Even/Odd positive whole number}

\begin{lstlisting}[rulesepcolor=\color{blue}]
writeInfo: ""
form
    natural number
endform
if number mod 2 = 0
    appendInfoLine: "The number ", number, " is even."
else
    appendInfoLine: "The number ", number, " is odd."
endif
\end{lstlisting}
\vspace{10pt}
% \\\\\\\\\\\\\\\\\\\\\\\\\\\\\\\\\\\\\\\\\\\\\\\\\\\\\\\\\\

% \\\\\\\\\\\\\\\\\\\\\\\\\\\\\\\\\\\\\\\\\\\\\\\\\\\\\\\\\\
\subsection*{\emph{9.} Make tones from fundamental frequency}

\begin{lstlisting}
writeInfoLine: ""
form
	natural tones
endform
fzero = 100
num = 0
for n 1 to tones
	num = num + 1
	Create Sound as pure tone: "tone", 1, 0, 0.4, 44100, fzero * num, 0.2, 0.01, 0.01
	Play
	appendInfoLine: fzero * num
	Remove
endfor
\end{lstlisting}

\begin{lstlisting}[rulesepcolor=\color{blue}]
form: "Play tones"
	natural number
endform
writeInfo: ""
for i to number
	frequency = i * 100
	Create Sound as pure tone: "tone", 1, 0, 0.4, 44100, frequency, 0.2, 0.01, 0.01
Play
appendInfoLine: "The frequency of sound ", i, " is ", frequency, " Hz."
Remove
endfor
\end{lstlisting}

\vspace{10pt}
% \\\\\\\\\\\\\\\\\\\\\\\\\\\\\\\\\\\\\\\\\\\\\\\\\\\\\\\\\\

% \\\\\\\\\\\\\\\\\\\\\\\\\\\\\\\\\\\\\\\\\\\\\\\\\\\\\\\\\\
\subsection*{\emph{10.} Make odd multiple tones from fundamental frequency}

\begin{lstlisting}
writeInfoLine: ""
form
	real fzero
	natural tones
endform
num = 0
for n 1 to tones
	num = num + 1
	odd = num * 2 - 1
	Create Sound as pure tone: "tone", 1, 0, 0.4, 44100, fzero * odd, 0.2, 0.01, 0.01
	Play
	appendInfoLine: fzero * odd
	Remove
endfor
\end{lstlisting}

\begin{lstlisting}[rulesepcolor=\color{blue}]
form: "Play odd overtones"
	real fzero
	natural tones
endform
writeInfo: ""
for n to tones
	freq = (2 * n - 1) * fzero
	Create Sound as pure tone: "tone", 1, 0, 0.4, 44100, freq, 0.2, 0.01, 0.01
	Play
	appendInfoLine: "The frequency of sound ", n, " is ", freq, " Hz."
	Remove
endfor
\end{lstlisting}

\vspace{10pt}
% \\\\\\\\\\\\\\\\\\\\\\\\\\\\\\\\\\\\\\\\\\\\\\\\\\\\\\\\\\

% \\\\\\\\\\\\\\\\\\\\\\\\\\\\\\\\\\\\\\\\\\\\\\\\\\\\\\\\\\
\subsection*{\emph{11.} Rolling a die until six}

\begin{lstlisting}[rulesepcolor=\color{blue}]
roll = 0
repeat
	die = randomInteger (1, 6)
	roll = roll + 1
until die = 6
writeInfoLine: "It took ", roll, " roll(s) for a six."
\end{lstlisting}
\vspace{10pt}
% \\\\\\\\\\\\\\\\\\\\\\\\\\\\\\\\\\\\\\\\\\\\\\\\\\\\\\\\\\


% \\\\\\\\\\\\\\\\\\\\\\\\\\\\\\\\\\\\\\\\\\\\\\\\\\\\\\\\\\
\subsection*{\emph{12.} Is praat-generated die fair?}

\begin{lstlisting}
turn = 0
count = 0
while turn < 10000
	die = randomInteger (1, 6)
	if die = 6
		count = count + 1
	endif
	turn = turn + 1
endwhile
writeInfoLine: "A six was rolled ", count, " times."
appendInfoLine: "Probability of a six roll: ", count / 10000
appendInfoLine: "Expected: 0.1666 (1/6) times, ie. around 1667 six rolls."
\end{lstlisting}

\begin{lstlisting}[title=\textbf{Info Window}]
A six was rolled 1676 times.
Probability of a six roll: 0.1676
Expected: 0.1666 (1/6) times, ie. around 1667 six rolls.
\end{lstlisting}

\begin{lstlisting}[rulesepcolor=\color{blue}]
count = 0
for turn to 10000
    die = randomInteger (1, 6)
    if die = 6
        count += 1
    endif
endfor
writeInfoLine: "The number of sixes rolled was ", count, "."
appendInfoLine: "I had expected ", 10000/6, " sixes."
\end{lstlisting}

\vspace{10pt}

Testing fairness of a die in this way is problematic as only the chance of a six getting rolled is observed for the 10000 turns, and not the same for all sides of the die. For the die to be fair, all sides need to have 1/6 chance. 
\vspace{10pt}
% \\\\\\\\\\\\\\\\\\\\\\\\\\\\\\\\\\\\\\\\\\\\\\\\\\\\\\\\\\

% ><><><><><><><><><><><><><><><><><><><><><><><><><><><><><
\subsection*{\emph{13.} The specific research question of your experiment}
\vspace{10pt}

\begin{itemize}
    \item \textbf{General research question:} \par
    Location of phoneme boundary as a function of orthographic context (and exopure/familiarity with orthography of certain language(s))
    \item \textbf{Addition:} \par
    L2 lexical knowledge context
    \item \textbf{Specific research question:}
    My idea is to take both orthography and non-native lexical knowledge as a coupled effect on phonological categorisation. Preliminary thoughts on operationalising this question: looking at phonemic boundary between voiced-plain-aspirated, aspirated-fricative continua for L2 English speakers from Assam (or L1 Assamese speaking). Onset stop contrast on L2 variety is through voicing unlike aspiration in L1 English for Assamese (voicing+aspiration for Boro), Orthograhpic <b>, <d>, <g> are percieved as voiced like other Indian Englishes. This perception gets carried over even when exposed to L1 English. Questions: How much closure devoicing / afterburst duration can be blocked off during phoneme categorisation for a voicing / plain percept (how strong would lexical effect be on a word like 'ban' / 'pan', and also orthographic; eg: when shown <pan> / <phan or fan>)? 
\end{itemize}
% ><><><><><><><><><><><><><><><><><><><><><><><><><><><><><

%wwwwwwwwwwwwwwwwwwwwwwwwwwwwwwwwwwwwwwwwwwwwwwwwwwwwwwwww
\end{document}